% page 379 du fac-simile (image p-378 du scan)
% bloc 165.1 x 242.0 mm ; marge gauche 36.9 mm
\pgc{15.240mm}{0.000mm}{
\l{\cel{55}371}
\l{}
\l{emezotinto.Graburo por qua onu preparas, unesme, la grano sur}
\l{\cel{3}la tota planko, til obtenar nigro perfekta, ed ube, pos de-}
\l{\cel{3}kkalqui ll ddesegnur ssu ll kkupro-plako oon aaplasta ppl o}
\l{\cel{4}min multe la grano per skrapilo por la loki klara, la mi-}
\l{\cel{3}nnuanci ll mmi-ombri - DDEFI}
\l{}
\l{emezurar. (\sou{trans., per}). Evaluar (quanto)per lua raporto ye quan-}
\l{\cel{4}to determinita, sama-speca, selektita kom komparo-termino.-}
\l{\cel{4}EFI.}
\l{}
\l{\cel{1}\sou{\textquotedbl{}mi\textquotedbl{}.}\cel{2}(\sou{muziko}). Noto triesma di la gamo diatonika naturala,}
\l{\cel{4}t.e. di la \textquotedbl{}do\textquotedbl{}-gamo.}
\l{}
\l{\cel{1}\sou{mi-}\cel{3}Prefixo di qua la senco esas \textquotedbl{}duime\textquotedbl{}.}
\l{}
\l{\cel{1}\sou{miasmo}. Emanajo des-salubra,nociva, de materii putranta. - DEFIR}
\l{}
\l{\cel{1}\sou{miaular}. (\sou{netrans.}) (\sou{aludante kato}). Agar la krio qua esas propr}
\l{\cel{4}ye lua speco. - DEFIS.}
\l{}
\l{\cel{1}\sou{micelo}. (\sou{biol.}) La aglomerajo maxim mikra ek molekuli, qua havas}
\l{\cel{4}la propraji fizikala di korpo.}
\l{}
\l{\cel{1}\sou{mielo}. Substanco siropatra, sukroza, quan la abeli elaboras ek}
\l{\cel{4}la materii quin li rekoliis en la flori, e quin li igas ek-}
\l{\cel{4}fluar aden la alveoli di lia abeluyo por la nutro di su e di}
\l{\cel{4}la larvi. - FIS.}
\l{}
\l{\cel{1}\sou{mielato}. Exuduro viskoza e sukroza, quan sekrecas, dum la perio-}
\l{\cel{4}di di aero-sikeso e di tero-sikeso, la folii di ula arbori.}
\l{}
\l{\cel{1}\sou{mieno}. Aspekto di la vizajo. - DEFR.}
\l{}
\l{\cel{1}\sou{migrar}. (\sou{netrans., ad, de)}. Diplasar su, livante ula lando por}
\l{\cel{4}instalar su en altra. - Su-diplasar di ula sorti de animali}
\l{\cel{4}qui iras ad altra klimati, sive periodale, segun la sezoni,}
\l{\cel{4}sive pro cirkonstanci acidentala. - EFIS.}
\l{}
\l{\cel{1}migreno. Doloro qua afektas, ordinare, nur parto di la kapo,}
\l{\cel{4}partikulare la regiono temporala, ed akompanesas multa-kaze}
\l{\cel{4}da perturbesi stomakala. - DEFIR.}
\l{}
\l{\cel{1}mikao. (\sou{mineral.}) Silikato aluminoza, kun metal-brilo, quan}
\l{\cel{4}onu povas separar facile ad lami diafana, per klivo.- DEFIS.}
\l{}
\l{\cel{1}mikologo.\cel{2}Autoro di traktato pri la fungi (champinioni).}
\l{}
\l{\cel{1}n\cel{2}- -- -. \textquotedbl{}- ---+-\cel{5}-+--\cel{2}- --- -+-\cel{2}--\cel{2}-\cel{2}---\cel{2};- ---,-}
\l{\cel{4}nionl).}
\l{}
\l{\cel{1}s\cel{1}\sou{korizo.}\cel{1}Fungo quan onu renkontras generale asociita kun radi-}
\l{\cel{4}Ej di vejetanti (precipue la radiki di la vejetanti montala;}
\l{\cel{4}kestanieri, avelanieri, fagi, morusieri. monto-pini, e o ) -}
\l{\cel{4}D ylS.}
}
